\section{The same native blowup across affine gluing constructions}

\begin{lemma}[The original affine Rees projective schemes agree]
\label{lem:fc-hironaka-dejong-affine-rees-comparison}
\leanok
\lean{EnvelopeProperty.HironakaDeJongReesComparison.grading_eq, EnvelopeProperty.HironakaDeJongReesComparison.toDeJong, EnvelopeProperty.HironakaDeJongReesComparison.toHironaka, EnvelopeProperty.HironakaDeJongReesComparison.affineIso, EnvelopeProperty.HironakaDeJongReesComparison.affineIso_structure}
The grading of the actual Rees subalgebra defined by vanishing of all
other polynomial coefficients is the grading defined by monomials of
the specified degree.  The identity of the original Rees subalgebra
therefore induces an isomorphism of their entire native projective
schemes, over the same original coefficient spectrum.
\end{lemma}
\begin{proof}
A polynomial whose coefficients outside degree $d$ vanish is exactly
the monomial formed from its coefficient in degree $d$.  Conversely
such a monomial has every other coefficient zero.  This proves
equality of every graded submodule.  The identity ring maps in both
directions are consequently graded homomorphisms and are surjective;
the irrelevant-ideal conditions are derived from surjectivity.  The
induced native projective maps are inverse, by functoriality of
projective spectrum and the two literal identity composites.  In
degree zero both maps preserve the original coefficient as the same
constant polynomial.  Naturality of the degree-zero projective
structure morphism gives equality of the original projections.
\end{proof}

\begin{lemma}[Compatibility with every original coefficient restriction]
\label{lem:fc-hironaka-dejong-rees-restriction}
\leanok
\lean{EnvelopeProperty.HironakaDeJongReesComparison.graded_naturality, EnvelopeProperty.HironakaDeJongReesComparison.affineIso_naturality}
\uses{lem:fc-hironaka-dejong-affine-rees-comparison}
For an original ring homomorphism $R\to S$ carrying an ideal to its
extension, the preceding entire affine-blowup isomorphism commutes
with the two actual original higher-level Rees morphisms and their
original coefficient maps.
\end{lemma}
\begin{proof}
Both homomorphisms on the original Rees subalgebras apply the same
ring homomorphism to every coefficient of the same polynomial.
Composing with either identity comparison therefore gives the same
graded ring homomorphism.  Contravariant functoriality of native
projective spectrum gives the asserted equality of entire scheme
morphisms.  Their original coefficient squares agree by the
preceding degree-zero comparison.
\end{proof}

\begin{theorem}[The whole original ideal-sheaf blowup comparison]
\label{thm:fc-hironaka-dejong-global-blowup-comparison}
\leanok
\lean{EnvelopeProperty.HironakaDeJongBlowupComparison.comparison, EnvelopeProperty.HironakaDeJongBlowupComparison.chart_comparison, EnvelopeProperty.HironakaDeJongBlowupComparison.comparison_projection, EnvelopeProperty.HironakaDeJongBlowupComparison.isPullback_chart_comparison, EnvelopeProperty.HironakaDeJongBlowupComparison.comparison_isIso, EnvelopeProperty.HironakaDeJongBlowupComparison.globalIso, EnvelopeProperty.HironakaDeJongBlowupComparison.globalIso_structure}
\uses{lem:fc-hironaka-dejong-rees-restriction}
For an ideal sheaf $I$ on an arbitrary original scheme $X$, the
blowup glued over all affine opens and the blowup glued over the
small affine Zariski site are isomorphic over the same original $X$.
\end{theorem}
\begin{proof}
Every object of the small affine site is an original affine open.
On it use the preceding entire affine-blowup comparison followed by
the canonical inclusion in the blowup glued over all affine opens.
The restriction comparison makes these maps a cocone for the small
site diagram, hence gives a morphism of the two whole glued schemes.
The component coefficient squares give its original projection
triangle.  Over every original affine open of $X$, both gluing
constructions are the pullback of their whole projection, and the
constructed morphism is the preceding affine isomorphism under
these two pullback identifications.  Thus it is an isomorphism on the
inverse images of the original affine cover.  Being an isomorphism
is local on the target, so the whole morphism is an isomorphism over
$X$, with the stated literal original projection equality.
\end{proof}
